%% =============================================================
%% Section 5: Conclusion
%% =============================================================

\section{Conclusion}
\label{sec:conclusion}

This paper develops a two-period small open economy New Keynesian model to compare
explicit forward guidance with opacity when the central bank's inflation aversion is
private information. The main findings are:

\begin{enumerate}[label=(\roman*)]
  \item Forward guidance weakly dominates opacity in welfare terms. The gain is
        strictly positive whenever there is genuine type uncertainty
        ($p \in (0,1)$, $\lambda_H \neq \lambda_L$) and cost-push shocks are present
        ($\sigma_c^2 > 0$).

  \item The welfare gain operates through two channels: an exchange rate channel---
        specific to the open-economy setting---and an expectation coordination channel.
        Both channels reduce the misalignment between private-sector beliefs and the
        central bank's actual reaction function.

  \item The value of transparency is non-monotone in the degree of openness: greater
        openness amplifies policy transmission but also compresses the cross-type
        dispersion in reaction functions, potentially reducing the informational value
        of the announcement.
\end{enumerate}

[Conclusion will be extended once quantitative results are available.]
